\documentclass[10pt,a4paper]{amsart}
\usepackage[margin=19mm]{geometry}
\usepackage{amsmath,amssymb,mathtools}
\usepackage[scr=rsfs,cal=euler]{mathalfa}
\usepackage{xfrac}
\usepackage[unicode,hidelinks,bookmarks=false]{hyperref}
\newcommand{\ie}{i.e.\ }
\newcommand{\cf}{cf.\ }
\newcommand{\resp}{respectively\ }
\newcommand{\Subsection}{\subsection}
\providecommand{\nobreakdash}{\nobreak\hbox{-}}
\setlength{\emergencystretch}{2em}
\multlinegap0pt
\usepackage{amssymb}
\usepackage{tikz}
\newtheorem{theorem}{Theorem}
\newcommand{\comp}{\mathrm{comp}}
\newcommand{\rev}{\mathrm{rev}}
\title{\textbf{Optimal Equivariant Matchings on the 6-Cube}\\[0.3em] \large With an Application to the King Wen Sequence}
\author{Alejandro Radisic}
\date{Source: arXiv:2601.07175v3 (CC BY 4.0)}
\begin{document}
\maketitle

\noindent\emph{Source and licence.} \textbf{Optimal Equivariant Matchings on the 6-Cube}\\[0.3em] \large With an Application to the King Wen Sequence. Version arXiv:2601.07175v3 (CC BY 4.0), distributed by arXiv under the Creative Commons Attribution 4.0 International licence, http://creativecommons.org/licenses/by/4.0/, as recorded in the arXiv licence metadata for this exact version and shown on the abstract page https://arxiv.org/abs/2601.07175v3. Full licence notice: \texttt{licenses/arxiv-2601.07175-CC-BY-4.0.txt}.\par\smallskip

\noindent\textit{Editorial scope.} Counted unit: the article's theorem thm:involution (source0.tex lines 331--333). The article's complete original proof of it is delivered with it (source0.tex lines 335--347, one \texttt{proof} environment of 13 lines). No mathematics is written, repaired or re-derived: the statement and the proof are the article's own lines, in the article's own notation, and no retained line is substituted or re-flowed.  No further source-local context is needed: every cross-reference of every retained line resolves inside the two delivered spans.  Nothing was omitted from any retained span, so the package carries no omission record.  No retained line contains a citation, so the wrapper carries no bibliography and no bibliography entry.  No retained line contains a figure environment, an image-inclusion command or drawing code, so no image is delivered and the package declares no asset dependency.  Editorial formatting only, in addition: an AMS \texttt{amsart} preamble that restates the article's own declarations and macros used by the retained lines, namely the environments \texttt{theorem} on separate counters, together with the standard packages those lines need.  The retained environments print in this document as Theorem 1 in order of appearance, whereas the article prints the same items with the numbering of its own full document.  The complete native source of the article as distributed in the arXiv e-print for this exact version is bundled unchanged as \texttt{sources/192483/source0.tex}.
\begin{theorem}[Involution]\label{thm:involution}
The priority partner function is an involution: $\mathrm{partner}(\mathrm{partner}(h)) = h$.
\end{theorem}

\begin{proof}
Two cases:
\begin{enumerate}
    \item If $h$ is a palindrome, then $\comp(h)$ is also a palindrome, so
    \[
    \mathrm{partner}(\comp(h))=\comp(\comp(h))=h.
    \]
    \item If $h$ is not a palindrome, then $\rev(h)$ is also not a palindrome, so
    \[
    \mathrm{partner}(\rev(h))=\rev(\rev(h))=h.
    \]
\end{enumerate}
\end{proof}

\end{document}
